\documentclass[11pt]{article}

\usepackage[margin=1in]{geometry}
\usepackage{amsmath,amssymb,amsthm}
\usepackage{booktabs}
\usepackage{hyperref}

\newtheorem{theorem}{Theorem}
\newtheorem{corollary}[theorem]{Corollary}
\newtheorem{remark}[theorem]{Remark}

\title{Why an Inverse-Elliptical Potential Produces a\\
       Force-Magnitude Dip Along its Major Axis}
\author{}
\date{}

\begin{document}
\maketitle

\begin{abstract}
For the artificial potential field $U=1/\rho$ built on an anisotropic
(elliptical) pseudo-distance $\rho$ with elongation ratio $q$, the force
magnitude $|\nabla U|$ has a strict local \emph{minimum} on the major axis
whenever $q>\sqrt{3}$. We derive the exact threshold, the location of the
off-axis ridge, and the depth of the dip in closed form, and we give the
geometric reason for the constant $\sqrt{3}$. This explains the narrow
low-force channel observed directly ahead of an obstacle in the DIPF force
plots, and shows that the feature is a necessary consequence of elongating
the field rather than a numerical artefact.
\end{abstract}

\section{Setup}

Place the obstacle at the origin and define the \emph{elliptical radius}
\begin{equation}
  \rho(x,y)\;=\;\sqrt{\,x^{2}+\frac{y^{2}}{q^{2}}\,},
  \qquad q>0 ,
  \label{eq:rho}
\end{equation}
where $q$ is the elongation of the field along $y$. In the implementation
$q=t_{y}/t_{x}$, the ratio of the two velocity-inflated time constants.
The potential is the inverse of this pseudo-distance,
\begin{equation}
  U(x,y)\;=\;\frac{1}{\rho(x,y)} .
  \label{eq:U}
\end{equation}
The level sets $\{U=\text{const}\}$ are concentric ellipses with semi-axis
ratio $q:1$, the major axis lying along $y$. The force is $F=-\nabla U$, and
the quantity plotted as the background colour field is $|F|=|\nabla U|$.

Differentiating \eqref{eq:U},
\begin{equation}
  \partial_{x}U=-\,x\,\rho^{-3},
  \qquad
  \partial_{y}U=-\,\frac{y}{q^{2}}\,\rho^{-3},
  \label{eq:partials}
\end{equation}
so that
\begin{equation}
  |\nabla U|^{2}
  \;=\;\frac{x^{2}+y^{2}/q^{4}}{\rho^{6}}
  \;=\;\frac{x^{2}+y^{2}/q^{4}}{\bigl(x^{2}+y^{2}/q^{2}\bigr)^{3}} .
  \label{eq:gradsq}
\end{equation}

\section{The dip criterion}

Fix a height $y=s>0$ and scan horizontally. Introduce the \emph{core
half-width} at that height,
\begin{equation}
  a\;=\;\frac{s}{q},
\end{equation}
which is the semi-minor axis of the equipotential ellipse passing through
$(0,s)$. Substituting $y=s$ into \eqref{eq:gradsq} and changing variables to
$u=x^{2}\ge 0$ gives
\begin{equation}
  g(u)\;:=\;\bigl|\nabla U\bigr|^{2}\Big|_{y=s}
  \;=\;\frac{u+a^{2}/q^{2}}{(u+a^{2})^{3}} .
  \label{eq:g}
\end{equation}
Because $|\nabla U|^{2}$ depends on $x$ only through $u=x^{2}$, it is an even
function of $x$ and its $x$-derivative vanishes at $x=0$ automatically. The
substantive question is therefore not whether $x=0$ is a critical point, but
whether it is a minimum or a maximum. That is decided by the sign of $g'$.

\begin{theorem}[Dip criterion]
\label{thm:main}
Let $g$ be as in \eqref{eq:g}. Then
\begin{equation}
  g'(u)\;=\;\frac{a^{2}\!\left(1-\dfrac{3}{q^{2}}\right)-2u}
                 {(u+a^{2})^{4}} ,
  \qquad u\ge 0 .
  \label{eq:gprime}
\end{equation}
Consequently:
\begin{enumerate}
  \item If $q\le\sqrt{3}$, then $g'(u)\le 0$ for all $u\ge0$, so $|\nabla U|$
        is non-increasing in $|x|$ and the major axis carries the
        \emph{largest} force on the scan line. There is no dip.
  \item If $q>\sqrt{3}$, then $g'(0)>0$, so $x=0$ is a strict local
        \emph{minimum} of $|\nabla U|$, and the unique maximum on the scan
        line lies at
        \begin{equation}
          x^{*}\;=\;\frac{a}{\sqrt{2}}\sqrt{1-\frac{3}{q^{2}}}
          \;=\;\frac{s}{q\sqrt{2}}\sqrt{1-\frac{3}{q^{2}}}
          \;\xrightarrow[\;q\to\infty\;]{}\;\frac{s}{q\sqrt{2}} .
          \label{eq:xstar}
        \end{equation}
\end{enumerate}
\end{theorem}

\begin{proof}
Write $g(u)=N(u)/D(u)$ with $N=u+a^{2}/q^{2}$ and $D=(u+a^{2})^{3}$. Then
$N'=1$ and $D'=3(u+a^{2})^{2}$, so by the quotient rule
\begin{align}
  g'(u)
  &=\frac{(u+a^{2})^{3}-\bigl(u+a^{2}/q^{2}\bigr)\cdot 3(u+a^{2})^{2}}
         {(u+a^{2})^{6}} \notag\\[2pt]
  &=\frac{(u+a^{2})-3\bigl(u+a^{2}/q^{2}\bigr)}{(u+a^{2})^{4}}
   \;=\;\frac{a^{2}\!\left(1-\dfrac{3}{q^{2}}\right)-2u}{(u+a^{2})^{4}} ,
\end{align}
where the second line cancels one factor of $(u+a^{2})^{2}$. This is
\eqref{eq:gprime}.

Since $a>0$, the denominator $(u+a^{2})^{4}$ is strictly positive on
$u\ge 0$. Hence $\operatorname{sign} g'(u)=\operatorname{sign}\bigl[
a^{2}(1-3/q^{2})-2u\bigr]$, an affine and strictly decreasing function of $u$.

If $q^{2}\le 3$ the constant term $a^{2}(1-3/q^{2})$ is non-positive, so the
numerator is $\le 0$ throughout $u\ge0$; thus $g$ is non-increasing, and so is
$|\nabla U|=\sqrt{g}$ (the square root is monotone). Since $u=x^{2}$ increases
with $|x|$, claim~1 follows.

If $q^{2}>3$ the numerator is positive at $u=0$ and decreases affinely,
vanishing exactly once at
\begin{equation}
  u^{*}=\tfrac{1}{2}a^{2}\!\left(1-\tfrac{3}{q^{2}}\right)>0 ,
  \label{eq:ustar}
\end{equation}
after which it is negative. Hence $g$ strictly increases on $[0,u^{*}]$ and
strictly decreases on $[u^{*},\infty)$: $u=0$ is a strict local minimum and
$u^{*}$ is the unique global maximum. Taking $x^{*}=\sqrt{u^{*}}$ gives
\eqref{eq:xstar}.
\end{proof}

\begin{remark}
A circle ($q=1$) can never produce the dip, and any elongation beyond
$q=\sqrt{3}\approx 1.732$ must. The feature is thus not an artefact of the
particular potential used, but a property of elongation itself.
\end{remark}

\section{Depth of the dip}

\begin{theorem}[Contrast ratio]
\label{thm:contrast}
For $q>\sqrt3$, the ratio of the ridge force to the on-axis force on any
horizontal scan line is
\begin{equation}
  C\;:=\;\frac{|\nabla U|_{\max}}{|\nabla U|_{x=0}}
  \;=\;\frac{2q^{3}}{3\sqrt{3}\,\bigl(q^{2}-1\bigr)}
  \;\xrightarrow[\;q\to\infty\;]{}\;\frac{2q}{3\sqrt{3}}\;\approx\;0.385\,q .
  \label{eq:contrast}
\end{equation}
In particular $C$ is independent of the height $s$.
\end{theorem}

\begin{proof}
From \eqref{eq:g}, on the axis
\begin{equation}
  g(0)=\frac{a^{2}/q^{2}}{a^{6}}=\frac{1}{q^{2}a^{4}} .
\end{equation}
At the ridge, substitute $u^{*}$ from \eqref{eq:ustar}. The two brackets
simplify because the $3/q^{2}$ terms combine:
\begin{align}
  u^{*}+\frac{a^{2}}{q^{2}}
  &=\frac{a^{2}}{2}\left(1-\frac{3}{q^{2}}\right)+\frac{a^{2}}{q^{2}}
   =\frac{a^{2}}{2}\left(1-\frac{1}{q^{2}}\right), \\[4pt]
  u^{*}+a^{2}
  &=a^{2}\left[1+\frac{1}{2}\left(1-\frac{3}{q^{2}}\right)\right]
   =\frac{3a^{2}}{2}\left(1-\frac{1}{q^{2}}\right).
\end{align}
Writing $\beta=1-q^{-2}$,
\begin{equation}
  g(u^{*})
  =\frac{\tfrac{1}{2}a^{2}\beta}{\bigl(\tfrac{3}{2}a^{2}\beta\bigr)^{3}}
  =\frac{\tfrac{1}{2}a^{2}\beta}{\tfrac{27}{8}a^{6}\beta^{3}}
  =\frac{4}{27\,a^{4}\beta^{2}} .
\end{equation}
Therefore
\begin{equation}
  C^{2}=\frac{g(u^{*})}{g(0)}
  =\frac{4}{27\,a^{4}\beta^{2}}\cdot q^{2}a^{4}
  =\frac{4q^{2}}{27\,\beta^{2}} ,
\end{equation}
and taking the positive square root with $\beta=(q^{2}-1)/q^{2}$ gives
\begin{equation}
  C=\frac{2q}{3\sqrt3\,\beta}=\frac{2q^{3}}{3\sqrt3\,(q^{2}-1)} .
\end{equation}
Both $a$ and $s$ have cancelled, proving the height-independence.
\end{proof}

\begin{corollary}[Geometry of the channel]
Since $C$ does not depend on $s$ while $x^{*}\propto s$ by \eqref{eq:xstar},
the locus of ridge points is a pair of straight rays through the obstacle,
\begin{equation}
  x=\pm\,\kappa\,s,
  \qquad
  \kappa=\frac{1}{q}\sqrt{\frac{1}{2}\left(1-\frac{3}{q^{2}}\right)} ,
\end{equation}
enclosing a wedge of uniformly suppressed force. The dip therefore appears in
a plot as a narrow channel of constant contrast opening away from the
obstacle, i.e.\ the observed ``U'' shape, rather than as an isolated blob.
\end{corollary}

\section{Why the constant is \texorpdfstring{$\sqrt{3}$}{sqrt(3)}}

The threshold has a direct geometric reading. Factor the gradient through the
metric:
\begin{equation}
  |\nabla U|
  =\bigl|U'(\rho)\bigr|\;\bigl|\nabla\rho\bigr|
  =\underbrace{\rho^{-2}}_{\text{amplitude}}\;
   \underbrace{\bigl|\nabla\rho\bigr|}_{\text{metric}},
  \qquad
  \bigl|\nabla\rho\bigr|=\frac{\sqrt{x^{2}+y^{2}/q^{4}}}{\rho} .
\end{equation}
Moving off-axis at fixed height, the two factors oppose one another:
\begin{itemize}
  \item $|\nabla\rho|$ \textbf{increases} from $1/q$ on the major axis to $1$
        on the minor axis. Equipotentials are $q$ times more widely spaced at
        the poles of the ellipse than at its flanks, so the force there is
        $q$ times weaker. This is the elongation, and it favours an off-axis
        ridge.
  \item $\rho^{-2}$ \textbf{decreases}, since $\rho$ grows with $|x|$. This is
        the inverse-distance decay, and it favours a monotone falloff.
\end{itemize}
The exponent $3$ appearing in \eqref{eq:gprime} is precisely the power
$\rho^{-3}$ carried by the gradient of a $1/\rho$ potential, as seen in
\eqref{eq:partials}. The dip survives only when the metric gain outpaces that
decay, which requires $q^{2}>3$. A different radial power law shifts the
threshold accordingly; an isotropic field has no metric gain at all
($|\nabla\rho|\equiv1$) and so never exhibits the dip.

\section{Numerical verification}

Evaluating $|\nabla U|$ by central differences on \eqref{eq:U} at $s=8$ over a
fine grid in $x$ reproduces both closed forms to four decimal places.

\begin{center}
\begin{tabular}{rcccc}
\toprule
$q$ & $x^{*}$ predicted & $x^{*}$ numerical & $C$ predicted & $C$ numerical\\
\midrule
$1.50$        & ---      & ---      & $1.0000$  & $1.0000$\\
$\sqrt{3}$    & $0.0000$ & $0.0149$ & $1.0000$  & $1.0000$\\
$2$           & $1.4142$ & $1.4142$ & $1.0264$  & $1.0264$\\
$3$           & $1.5396$ & $1.5396$ & $1.2990$  & $1.2990$\\
$6$           & $0.9027$ & $0.9027$ & $2.3754$  & $2.3754$\\
$12$          & $0.4665$ & $0.4665$ & $4.6511$  & $4.6511$\\
$30$          & $0.1882$ & $0.1883$ & $11.5598$ & $11.5598$\\
\bottomrule
\end{tabular}
\end{center}

\noindent
At $q=\sqrt3$ the maximum is degenerate ($x^{*}=0$) and the numerical value
merely reflects the grid resolution, consistent with $C=1$.

\section{Application to the implemented field}

The implemented potential carries an additional vehicle-dimension prefactor,
\begin{equation}
  U_{\text{impl}}(x,s)=\frac{\lambda}{\bigl(1-l/s\bigr)\,\rho(x,s)} ,
  \label{eq:Uimpl}
\end{equation}
with $l$ the obstacle length and $s$ the longitudinal gap. Two consequences:

\begin{itemize}
  \item \textbf{The ridge location is unchanged.} The prefactor
        $\bigl(1-l/s\bigr)^{-1}$ is constant in $x$ along a horizontal scan,
        so it rescales the $x$-profile without shifting its stationary point.
        For $q=12$ at $s=9.5$, \eqref{eq:xstar} predicts $x^{*}=0.554$ against
        a measured $0.54$ in the implementation.
  \item \textbf{The dip is shallower.} The factor does depend on $s$, so
        $\partial_{y}$ acts on it and contributes an additional,
        $x$-dependent term to $F_{y}$ that partially fills the on-axis
        minimum. The measured contrast is $C\approx2.52$ against the
        pure-elliptical value $4.6511$ --- qualitatively identical, roughly
        half as deep.
\end{itemize}

Finally, note that the dip lives in $|\nabla U|$, not in $U$. The potential
itself is maximal on the major axis and decreases monotonically outward, so
its level sets remain nested convex ellipses; it is the level sets of the
\emph{force magnitude} that become non-convex and pinched at the poles. In
control terms the major axis is a ridge of $U$, hence an unstable
equilibrium line: a host vehicle exactly on the centreline receives a
near-zero lateral command, while a small lateral offset receives a large one.

\end{document}
